\section{Anchored involutions and field descent}
The retained rotation correspondence uses generators
$\sigma_1,\ldots,\sigma_{n-1}$ and the interval-involution relations.
Its intersection axiom says that subgroups generated by interval products
with endpoints in $I$ and $J$ intersect in the subgroup for $I\cap J$.
We use this full axiom, not merely the rank-four special case.

\begin{lemma}\label{E.Anchored}
In the rotation group of a rank-$n$ chiral polytope, the involutions
$t_i=\sigma_i\cdots\sigma_{n-1}$, $1\le i\le n-2$, form a string
C-marking of rank $n-2$. Its adjacent labels are
$|\sigma_1|,\ldots,|\sigma_{n-3}|$, and its even subgroup
$\langle\sigma_1,\ldots,\sigma_{n-3}\rangle$ has index two.
There is a dual anchored marking with the labels at the other end.
\end{lemma}
\begin{proof}
Every $t_i$ is a nonidentity involution by the rotation correspondence.
For $i<j$, cancellation gives $t_it_j=\sigma_i\cdots\sigma_{j-1}$.
If $j-i\ge2$ this is an involution, proving the string commutations;
if $j=i+1$ it is $\sigma_i$. For a set of indices $I$, the group
$\langle t_i:i\in I\rangle$ is exactly the endpoint subgroup of the
rotation intersection axiom with vertex indices
$\{i-1:i\in I\}\cup\{n-1\}$. Products $t_it_j$ supply every
interval with both endpoints in this set. Intersecting two such sets
therefore proves the complete C-intersection property, including the
empty intersection. The even words in the $t_i$ are generated by the
adjacent products $\sigma_1,\ldots,\sigma_{n-3}$. This subgroup cannot
contain $t_{n-2}$: its endpoint set meets that for $t_{n-2}$ in a single
vertex, whose endpoint subgroup is trivial. Thus its index is two.
Reversing the indexing gives the dual statement.
\end{proof}

\begin{lemma}\label{E.NoTwoLabel}
If a chiral rotation group is $L(q)$ or $P(q)$, $q\ge4$, all its labels
are greater than two.
\end{lemma}
\begin{proof}
The first label cannot be two: the identity automorphism would then
satisfy the rotation reversal rule
$\sigma_1\mapsto\sigma_1^{-1}$,
$\sigma_2\mapsto\sigma_1^2\sigma_2$ and fix the other generators.
The same holds at the last label by duality.
Suppose instead $\sigma_j^2=1$ with $1<j<n-1$. Put
$A=\langle\sigma_1,\ldots,\sigma_{j-1}\rangle$ and
$B=\langle\sigma_{j+1},\ldots,\sigma_{n-1}\rangle$.
For $x=\sigma_i\cdots\sigma_{j-1}$ and
$y=\sigma_{j+1}\cdots\sigma_k$, all of $x\sigma_j$, $\sigma_j y$
and $x\sigma_j y$ are involutions. Thus $\sigma_j$ inverts each
such $x,y$, and the last relation gives $xy=yx$. These prefix and
suffix elements generate $A,B$. Hence $A$ and $B$ commute and are
normal in $\langle A,B,\sigma_j\rangle=G$; both are nontrivial.
This is impossible in $L(q)$ or $P(q)$: every nontrivial normal subgroup
contains the simple socle $L(q)$, whose centralizer is trivial.
These are exactly the standard structural inputs already retained in
Topic E. No assumption about a minimal generating number is used.
\end{proof}

\begin{lemma}\label{E.SubfieldEnvelope}
Let $k\subseteq K$ be finite fields of odd characteristic, $|k|\ge5$.
Any conjugate of $P(l)$ or $L(l)$ in $P(K)$ containing the standard
$L(k)$ has $k\subseteq l$ and is the standard $P(l)$ or $L(l)$.
Consequently its intersection with standard $P(k)$ is $P(k)$ in the
first case. In the second case it is $P(k)$ if $[l:k]$ is even and
$L(k)$ if $[l:k]$ is odd.
\end{lemma}
\begin{proof}
The projective invariant $J(X)=\operatorname{tr}(X)^2/\det(X)$ lies
in $l$ on such a conjugate. The matrices
$\left(\begin{smallmatrix}a&-1\\1&0\end{smallmatrix}\right)\in L(k)$
show that $a^2\in l$ for every $a\in k$. Squares span $k$, since
$a=((a+1)/2)^2-((a-1)/2)^2$, so $k\subseteq l$.

Let $\phi$ be the $|l|$-power field automorphism of $K$.
The standard $P(l)$ is its pointwise fixed projective group: normalize
one nonzero matrix entry to one to check this assertion directly.
A conjugate $P(l)^g$ is the fixed group of the conjugate semilinear
automorphism $\theta$. Both $\theta$ and $\phi$ fix $L(k)$ pointwise,
so their quotient is an inner projectivity centralizing $L(k)$.
That centralizer is trivial: it centralizes the upper and lower
nonidentity elementary unipotents, hence fixes infinity and zero,
and commutation with translation forces its diagonal ratio to be one.
Thus $\theta=\phi$ and $P(l)^g=P(l)$. The same argument applied to
an overgroup of $L(l)^g$ places it in standard $P(l)$; its simple
index-two subgroup is standard $L(l)$. Finally a scalar in $k^\times$
is a square in $l$ exactly when it is a square in $k$ or $[l:k]$ is
even. Applying this to determinants proves the intersections.
\end{proof}

We use the accepted Topic E classification: a projective group with a
rank-four C-marking is $P(5)$, $L(11)$ or $L(19)$, with its specified
marking. The proof and finite table belong to the complete common paper
and remain paid once. We also retain Dickson's proper-subgroup alternatives
in their natural-subfield form. Here are the precise additional interfaces.

\begin{lemma}\label{E.RegularSubgroups}
A directly regular rank-three C-group contained in $P(q)$, with both
labels greater than two, is either $S_4$ with its tetrahedral marking,
or a natural conjugate of $P(k)$ with $k$ odd, $|k|\ge5$, whose rotation
subgroup is $L(k)$. A directly regular rank-four C-group contained in
$P(q)$, with all labels greater than two, is the simplex $P(5)\cong S_5$.
No C-marking contained in $P(q)$ with these labels has rank at least five.
\end{lemma}
\begin{proof}
In odd characteristic an affine subgroup has only half-turn involutions;
two distinct such involutions cannot commute. It cannot contain a
rank-three string marking. In characteristic two its involutions are
translations and all commute. A cyclic group is likewise excluded.
In a dihedral group, commuting endpoint involutions of a rank-three
marking cannot be central, so are opposite reflections. Writing them
$s,sr^{N/2}$ and the middle reflection $sr^a$, generation says
$\gcd(a,N/2)=1$. At least one of $a,a-N/2$ is coprime to $N$; the
corresponding adjacent pair already generates the whole group,
contradicting C-intersection. A central middle involution would instead
give a label two.

All involutions in $A_4$ commute. The group $A_5$ has no index-two
subgroup. A directly regular marking of $S_4$ consists of transpositions,
since its even subgroup is its unique index-two subgroup $A_4$.
Commuting distinct endpoints are disjoint transpositions and a middle
transposition joining them gives precisely type $\{3,3\}$.
The remaining Dickson alternatives are natural projective subfield groups.
The simple $L(k)$ has no index-two subgroup; $P(k)$ has one exactly in
odd characteristic, namely $L(k)$. The small fields give the cases already
handled. This proves the rank-three assertion.

For rank four, the affine and dihedral exclusions still apply to the
first three generators. The involution centralizer in $A_5$ is a Klein
four group, so a first generator commuting with the noncommuting last
two is impossible. In $S_4$, if an endpoint were a transposition its
centralizer would be abelian, giving the same contradiction. Thus both
endpoints are double transpositions. The two middle generators must be
transpositions, since adjacent labels exceed two. The first adjacent
pair generates a dihedral group of order eight containing the entire
normal Klein four group, hence the last generator, contrary to the
intersection property. This also excludes $A_4$.
Dickson now reduces to projective subfield groups. The common rank-four
classification leaves $P(5),L(11),L(19)$; only $P(5)$ has an index-two
subgroup, and its marking is the simplex marking.

Finally, at rank at least five both fourth and fifth generators would
centralize the first noncommuting pair. The common projective centralizer
lemma allows at most one nonidentity element there in odd characteristic;
in characteristic two the centralizer of an involution is abelian and
already prevents rank four. This is independent of which subgroup the
marking generates.
\end{proof}

\begin{theorem}\label{E.ChiralRankBound}
Neither $L(q)$ nor $P(q)$, $q\ge4$, is the full group of a chiral
polytope of rank at least five.
\end{theorem}
\begin{proof}
All labels exceed two. Anchoring gives a directly regular C-subgroup of
rank $n-2$, so $n\le6$. If $n=6$, the anchored rank-four subgroup is a
simplex, so the first three labels are three. Its dual makes the last
three labels three. The full rotation group is therefore a quotient of
the even simplex Coxeter group $[3,3,3,3,3]^+\cong A_7$.
This group is simple, so the nontrivial quotient is itself, with its
regular simplex marking. This contradicts chirality.

For $n=5$ write the rotations as $\alpha,\beta,\gamma,\delta$.
The two anchored directly regular groups are
\[
 P_0=\langle\alpha,\beta,\gamma\delta\rangle,
 \qquad Q_0=\langle\alpha\beta,\gamma,\delta\rangle,
\]
with even subgroups $R=\langle\alpha,\beta\rangle$ and
$S=\langle\gamma,\delta\rangle$. If both are tetrahedral,
all four labels are three and the same argument uses the simple
$[3,3,3,3]^+\cong A_6$ and its regular marking.

First suppose the whole group is $P(q)$ with $q$ odd. The even group
of either anchored subgroup lies in $L(q)$: it is a natural $L(k)$,
or is $A_4$ generated by elements of order three, which must lie in
the determinant-square kernel. Thus all four rotations lie in $L(q)$,
a contradiction. If $q$ is even there are no odd-characteristic natural
$P(k)$ subgroups, so both anchored groups would be tetrahedral, already
excluded.

It remains that the whole group is $L(q)$ with $q$ odd, and at least
one anchored group, say $P_0$, is standard $P(k)$, $|k|\ge5$.
Then $R=L(k)$ and $[\mathbb F_q:k]$ is even. Put
$F_0=\langle\alpha,\beta,\gamma\rangle$, a proper subgroup by
rotation intersection. Dickson's alternatives containing $R$ are
projective subfield groups; the sole exceptional possibility is
$R=F_0=A_5$, which would imply
$G=\langle P_0,\gamma\rangle=P_0$, impossible.
The subfield-envelope lemma puts $F_0$ in standard $P(l)$ or $L(l)$
with $k\subseteq l$. The first possibility contains $P_0$, and again
$G=\langle F_0,P_0\rangle=F_0$ is impossible.
In the second, the exact rotation intersection is
$F_0\cap P_0=R=L(k)$. Thus $[l:k]$ is odd, while
$[\mathbb F_q:k]$ is even. It follows that
$[\mathbb F_q:l]$ is even, so $P(l)$ is a proper subgroup of $L(q)$.
But $G=\langle F_0,P_0\rangle\le P(l)$, the final contradiction.
\end{proof}
